\documentclass[11pt,a4paper]{article}
\usepackage[margin=1in]{geometry}
\usepackage[T2A]{fontenc}
\usepackage[utf8]{inputenc}
\usepackage[english,russian]{babel}
\usepackage{amsmath,amssymb}
\usepackage{microtype}
\usepackage{hyperref}

\setlength{\parindent}{0pt}
\setlength{\parskip}{0.75\baselineskip}
\sloppy

\title{Элементарная частица как замкнутая волна поля 01\\\large Скалярный профиль, интерпретация удержания нормали, масса, память и отсутствие точечной сингулярности}
\author{}
\date{}

\newcommand{\field}{\mathcal{F}_{01}}
\newcommand{\phase}{\varphi}
\newcommand{\normal}{\mathcal{N}}
\newcommand{\memory}{\mathcal{M}}
\newcommand{\contour}{\mathcal{C}}

\begin{document}
\maketitle

\begin{center}
\textit{Рабочая версия v0.2 --- русская исследовательская редакция}
\end{center}

\begin{abstract}
В статье формулируется первое ядро модели 01: элементарная частица рассматривается не как материальная точка, а как устойчивая замкнутая волновая конфигурация поля. Два режима поля --- фиксация $0$ и направление $1$ --- образуют ритм, из которого возникают фазовый обход, локальный скалярный/VEV-профиль, его интерпретация как удержания нормали, масса и память. Фотон интерпретируется как открытый обход фазы $1$, тогда как массивная элементарная частица возникает только при замыкании обхода и формировании локального скалярного/VEV-профиля, которому в модели 01 приписывается интерпретация удержания нормали. Такой подход позволяет заменить точечное ядро динамическим ядром и тем самым избежать внутренней сингулярности в описании элементарной частицы. Горизонт чёрной дыры появляется как предельный режим той же логики: подавление локального скалярного/VEV-профиля и переход фазовой памяти в поверхностную запись.
\end{abstract}

\textbf{Ключевые слова:} модель 01; замкнутая волна; элементарная частица; фазовый обход; скалярный профиль; VEV-профиль; интерпретация удержания нормали; масса; память; квантовая теория поля; горизонт; голография.

\section{Как читать эту статью}

Эта статья является рабочим исследовательским текстом, а не завершённой физической теорией. Её цель --- сделать первый шаг от книги и интуитивной модели к формальному языку, который можно обсуждать, критиковать и постепенно сопоставлять с известной физикой.

Читателю важно сразу различать три уровня утверждений:
\begin{enumerate}
\item \textbf{Известная физика}: квантовые поля, вакуумные корреляции, энергия фотона $E=\hbar\omega$, связь массы и энергии, тензор энергии--импульса, энтропия горизонта.
\item \textbf{Интерпретация модели 01}: частица описывается как устойчивый режим поля, фотон --- как открытый фазовый перенос, массивная частица --- как замкнутый узел с локальным скалярным/VEV-профилем, которому в модели 01 приписывается интерпретация удержания нормали.
\item \textbf{Открытая гипотеза}: скалярный/VEV-профиль, интерпретация удержания нормали, память, фазовое натяжение, поверхностная запись горизонта и отсутствие буквального испарения чёрной дыры в модели 01 требуют дальнейшей математической формализации.
\end{enumerate}

Поэтому все формулы, которые не выведены из строгого действия или из стандартной теории, следует читать как рабочие обозначения и направления будущей формализации. В частности, модель 01 не объявляет Стандартную модель неверной и не заменяет квантовую теорию поля. Она предлагает дополнительный фазово-топологический язык, который должен быть проверен на согласованность.

\section{Цель статьи}

Модель 01 предлагает рассматривать физические структуры не как набор готовых объектов в заранее данном пространстве, а как устойчивые режимы одного поля. В этой статье выделяется минимальное ядро модели: возможное описание элементарной частицы как замкнутой волны.

Главная задача статьи --- ввести рабочий словарь, который позволит перейти от образного описания к математическому каркасу. Мы не утверждаем, что данная формализация уже заменяет Стандартную модель. Напротив, она рассматривается как интерпретационный и гипотетический слой, который должен быть дальше сопоставлен с известными уравнениями квантовой теории поля.

Практическая цель версии v0.2 --- отделить физически установленные элементы от новых предположений модели. Это необходимо для того, чтобы дальнейшие статьи могли развивать не лозунг, а исследовательскую программу: определить динамические переменные, действие, энергетический функционал, устойчивые решения и возможные проверяемые отличия.

\section{Два режима поля 01}

Пусть $\field$ обозначает поле, имеющее два неразделимых режима:
\begin{align}
0 &: \text{фиксация, удержание, запись}, \\
1 &: \text{направление, движение, обход}.
\end{align}

Эти режимы не являются двумя отдельными веществами или независимыми полями. Они задают минимальное различие, из которого возникает ритм. Поэтому основная структура модели записывается не как сумма двух объектов, а как пара режимов:
\begin{equation}
\field = (0,1).
\end{equation}

Режим $0$ фиксирует различие, а режим $1$ разворачивает его. Без фиксации движение не оставляет следа; без направления фиксация не создаёт процесса. Поэтому физическая структура появляется только там, где возникает ритм между $0$ и $1$.

\section{Фаза и обход}

Для первого математического приближения введём фазу $\phase$, описывающую состояние ритма поля. Волновой режим можно записывать в обычной комплексной форме:
\begin{equation}
\psi = A e^{i\phase}.
\end{equation}

Однако в модели 01 фаза понимается не только как параметр волны, но и как мера различия между движением и фиксацией. Пока фаза свободно переносится без замыкания, возникает открытый режим. Когда фазовый обход замыкается, появляется устойчивый узел.

Условие замкнутого фазового обхода можно записать как
\begin{equation}
\oint_{\contour} d\phase = 2\pi n, \qquad n \in \mathbb{Z}.
\label{eq:phase_quantization}
\end{equation}

Это условие не вводится как окончательный закон частиц, а используется как минимальное выражение идеи: устойчивый узел возникает только тогда, когда обход возвращается к самосогласованному состоянию.

\section{Фотон как открытый обход}

Фотон в данной модели является открытым обходом фазы $1$. Он переносит изменение ритма, но не формирует локальный скалярный/VEV-профиль, которому в модели 01 приписывается удержание нормали. Поэтому фотон не имеет массы покоя.

Схематически:
\begin{equation}
\text{фотон}: \quad 1_{\mathrm{open}}, \qquad \normal = 0.
\end{equation}

Это означает не отсутствие физического состояния, а отсутствие локальной глубины. Фотон переносит фазовый отпечаток события: фазу, частоту, поляризацию и направление. Поэтому он может нести информацию о событии, которое его породило, но не является локальным узлом памяти.

В языке книги энергия фотона понимается не как вещество, помещённое внутрь частицы, а как интенсивность незамкнутого фазового переноса. Поэтому энергия фотона проявляется через частоту:
\begin{equation}
E=\hbar\omega.
\end{equation}
В модели 01 эта стандартная формула читается так: чем выше частота ритма, тем интенсивнее перенос фазы $1$.

\section{ЭЧ как замкнутая волна}

Элементарная частица в модели 01 определяется как устойчивая замкнутая интерференционная структура. Она не является материальной точкой. Её минимальная структура включает:
\begin{itemize}
\item две тангенциальные дуги обхода, связанные с режимом $1$;
\item локальный скалярный/VEV-профиль, которому в модели 01 соответствует интерпретация удержания нормали и связь с режимом $0$;
\item внутренний рисунок интерференции;
\item связь с внешним полем.
\end{itemize}

Схематически:
\begin{equation}
\text{ЭЧ}: \quad (1_{\tau_1},1_{\tau_2}) + \normal_0.
\end{equation}

Две тангенциальные дуги не являются двумя отдельными частями вещества. Это два фазовых направления, которые удерживают друг друга. В осторожной формализации их несовпадение сначала описывается локальным скалярным/VEV-профилем; язык нормали остаётся интерпретацией модели 01, а не самостоятельной новой степенью свободы.

\section{Скалярный профиль и интерпретация локальной глубины}

Первичным объектом здесь является локальный скалярный профиль $N$ или VEV-подобный профиль. Обозначение $\normal$ остаётся интерпретационным ярлыком модели 01: оно указывает на локальную глубину, связанную со сцеплением двух обходов.

Можно ввести профиль как функционал фазового разнесения:
\begin{equation}
N = N[\Delta\phase], \qquad \normal \equiv \text{интерпретация}(N),
\end{equation}
где $\Delta\phase$ обозначает несовпадение двух фазовых обходов.

Если фазовое разнесение исчезает,
\begin{equation}
\Delta\phase \to 0,
\end{equation}
то
\begin{equation}
N \to 0, \qquad \normal \to 0 \;\text{как интерпретационный предел}.
\end{equation}

В этом пределе локальная глубина исчезает: стандартный скалярный профиль подавляется, а замкнутый узел расплетается обратно в поле или переходит в поверхностный режим записи.

\section{Масса как энергия удержания скалярного профиля}

В модели масса не понимается как количество вещества внутри частицы. На уровне осторожной формализации масса связывается с энергией удержания локального скалярного/VEV-профиля; язык удержания нормали остаётся интерпретацией этого скалярного/VEV-слоя. Чем устойчивее поле удерживает локальную глубину, тем сильнее проявляется инерция.

В более общем смысле энергия в модели 01 не вводится как отдельная субстанция. Она описывает интенсивность фазовой динамики и способность поля менять, переносить или удерживать режим. Поэтому незамкнутый режим проявляется как перенос энергии, а замкнутый режим --- как энергия удержания структуры.

Схематически:
\begin{equation}
 m \sim \mu(N), \qquad \normal \equiv \text{интерпретация}(N),
\end{equation}
где $\mu$ --- мера устойчивости локального скалярного/VEV-профиля в данной интерпретации.

В качестве качественной аналогии можно записать массу как энергию эффективного натяжения вдоль замкнутого контура:
\begin{equation}
 m \sim \frac{1}{c^2}\oint_{\contour} T_{\mathrm{field}}\, dl.
\label{eq:mass_tension}
\end{equation}

Здесь $T_{\mathrm{field}}$ пока не является строго определённой компонентой тензора энергии--импульса. В данной версии статьи это эффективная величина, обозначающая локальную цену фазового натяжения вдоль замкнутой траектории. Формула \eqref{eq:mass_tension} является рабочим наброском и качественной аналогией, а не завершённым выводом. Её задача --- зафиксировать принцип: масса связана с энергией удержания замкнутой фазовой структуры.

В будущей формализации возможны два варианта. Либо $T_{\mathrm{field}}$ должен быть выведен из плотности энергии конкретного действия поля 01, либо он должен быть сопоставлен с проекцией тензора энергии--импульса на касательное направление контура:
\begin{equation}
T_{\mathrm{field}} \quad \longrightarrow \quad T_{\mu\nu}u^{\mu}u^{\nu}
\end{equation}
или с близкой по смыслу плотностью энергии устойчивой конфигурации. До такого вывода выражение \eqref{eq:mass_tension} следует читать только как указание направления формализации.

\section{Динамическое ядро вместо точечного ядра}

Если частица мыслится как точечный объект с внутренним материальным центром, возникает проблема сингулярности: вся структура оказывается сжатой в нулевой объём. В модели 01 это заменяется другим описанием.

У ЭЧ нет материального ядра в виде точки или плотного шарика. Но есть динамическое ядро:
\begin{equation}
\text{динамическое ядро} = \text{ось устойчивости} + \text{минимальный сердцевинный обход}.
\end{equation}

Такое ядро не является веществом внутри частицы. Это минимально самосогласованный режим движения. Спин в этом языке соответствует ориентации оси устойчивости, а не вращению твёрдого тела.

Отсюда следует важный вывод: память и масса ЭЧ не требуют точечной сингулярности. Они распределены в ритме, обходе и локальном скалярном/VEV-профиле, который в модели 01 получает интерпретацию удержания нормали.

\section{Память как межфазовая структура}

Память ЭЧ не принадлежит форме частицы как внешнему контуру. Она находится между фазами, то есть в волне, связывающей два обхода.

Обозначим память узла как $\memory$:
\begin{equation}
\memory = \memory(\Delta\phase, \contour, N; \normal_{\mathrm{interp}}).
\end{equation}

Эта запись означает, что память задаётся не одной координатой, а всей структурой фазового обхода: задержками, ориентациями, пересечениями и согласованиями.

В таком языке материя является не субстанцией, а устойчивым способом поля сохранять историю своих деформаций.

\section{Силы как режимы замкнутой волны}

После появления замкнутой волны возникают физические свойства, которые в обычном языке называются взаимодействиями. В модели 01 они интерпретируются как разные аспекты фазовой геометрии узла:
\begin{center}
\begin{tabular}{ll}
масса & энергия удержания локального скалярного/VEV-профиля, \\
заряд & ориентация фазового обхода, \\
спин & внутренняя структура фазового обхода, \\
электричество & асимметрия тангенциальных фаз, \\
магнетизм & вихревой режим движения фазы $1$, \\
гравитация & коллективная складка фазы $0$, \\
сильное взаимодействие & согласование локальных профилей соседних узлов, \\
слабое взаимодействие & перестройка или подавление локального профиля.
\end{tabular}
\end{center}

Сила в общем виде может рассматриваться как градиент фазового рельефа:
\begin{equation}
F \sim -\nabla \Phi(x),
\end{equation}
где $\Phi(x)$ --- эффективная фазовая геометрия поля.

\section{Связь с известной физикой}

Чтобы модель 01 не превращалась только в метафорический язык, важно сразу отделить два уровня описания. Первый уровень --- это установленная квантовая теория поля, где частицы являются возбуждениями полей, а само понятие частицы зависит от выбора мод, состояния наблюдателя и доступной области пространства-времени. Второй уровень --- интерпретационная гипотеза модели 01, где эти возбуждения описываются как открытые или замкнутые фазовые режимы поля.

В обычной квантовой теории поля уже есть важный нам мотив: вакуум не является пустотой. Даже если среднее поле равно нулю,
\begin{equation}
\langle 0|\phi(x)|0\rangle = 0,
\end{equation}
корреляционная функция может оставаться ненулевой:
\begin{equation}
G(x,y)=\langle 0|\phi(x)\phi(y)|0\rangle \neq 0.
\end{equation}
В языке модели 01 это можно читать осторожно: поле до появления локального узла уже обладает структурой корреляций. Элементарная частица тогда не создаётся из абсолютной пустоты, а возникает как устойчивое локальное замыкание уже существующей фазовой структуры.

Другой важный мотив --- зависимость частичного содержания от наблюдателя. В искривлённом пространстве-времени или в ускоренных системах разные разложения поля на моды могут быть связаны преобразованием Боголюбова:
\begin{equation}
a_{\mathrm{out}}=\alpha a_{\mathrm{in}}+\beta a_{\mathrm{in}}^{\dagger}.
\end{equation}
Если коэффициент $\beta$ ненулевой, один и тот же квантовый режим может давать разное эффективное описание частиц для разных наблюдателей. Для модели 01 это означает, что ``частица'' не должна пониматься как маленький шарик, существующий одинаково для всех описаний. Более естественно говорить о стабильном режиме поля, который становится частицей в той мере, в какой он даёт устойчивый отклик, массу, заряд, спин и канал взаимодействия.

Фотон в такой связке занимает особое место. В стандартной физике он является безмассовым квантом электромагнитного поля; в модели 01 это соответствует открытому фазовому переносу без удерживаемого локального скалярного/VEV-профиля. Массивный фермион, напротив, требует не только переноса фазы, но и устойчивой внутренней самосогласованности. Поэтому условие
\begin{equation}
\oint_{\contour} d\phase = 2\pi n
\end{equation}
следует понимать как топологический образ стабильности, а не как готовую замену уравнениям Дирака, Янга--Миллса или механизму Хиггса.

Механизм Хиггса в этой статье также не отменяется. В Стандартной модели массы элементарных частиц связаны с взаимодействием с полем Хиггса и калибровочной структурой теории. Модель 01 пока предлагает более глубокий интерпретационный вопрос: почему масса вообще проявляется как устойчивость локального режима? Ответ модели формулируется осторожнее: масса есть энергия удержания локального скалярного/VEV-профиля, которому затем приписывается интерпретация удержания нормали. Поэтому будущая формализация должна показать, можно ли связать параметр $\mu(N)$ с известными массовыми членами и константами связи, а не просто заменить их новым словом.

Спин также требует осторожности. В обычной квантовой теории спин не является вращением маленького тела; это внутреннее квантовое число, связанное с представлениями группы Лоренца. Поэтому фраза ``спин как структура обхода'' означает не механическое вращение, а фазово-топологический способ кодировать внутреннюю ориентацию узла. Особенно важным ориентиром здесь является фермионная периодичность: состояние со спином $1/2$ возвращается к себе только после поворота на $4\pi$, а не на $2\pi$. В языке модели 01 это может быть записано как условие скрученного обхода:
\begin{equation}
\Psi(\phase+2\pi)=-\Psi(\phase), \qquad \Psi(\phase+4\pi)=\Psi(\phase).
\end{equation}

Наконец, горизонт чёрной дыры связывает эту статью с информационной стороной квантовой теории поля. В стандартном описании горизонт ограничивает доступные степени свободы: если часть системы недоступна, полное состояние заменяется редуцированной матрицей плотности
\begin{equation}
\rho_{\mathrm{out}}=\operatorname{Tr}_{\mathrm{in}}\left(|\Psi\rangle\langle\Psi|\right).
\end{equation}
Это даёт энтропийное и термодинамическое описание горизонта. В модели 01 тот же мотив переформулируется как запись: объёмная фазовая структура замкнутой волны переходит в поверхностный код. Так появляется мост к голографическому принципу, где информация объёма выражается через границу.

Итак, связь с известной физикой пока имеет статус программы, а не завершённого вывода:
\begin{enumerate}
\item вакуумные корреляции дают язык для поля до локализации;
\item зависимость частиц от выбора мод поддерживает идею частицы как режима поля;
\item фазовая квантизация даёт топологический образ устойчивости;
\item масса как энергия удержания локального скалярного/VEV-профиля должна быть сопоставлена с механизмом Хиггса;
\item спин как скрученный обход должен быть сопоставлен с представлениями группы Лоренца;
\item горизонт как запись должен быть сопоставлен с редуцированными состояниями, энтропией и голографией.
\end{enumerate}

\section{Предел: подавление скалярного профиля и горизонт}

Если локальный скалярный/VEV-профиль стремится к нулю, то в интерпретации модели 01 исчезает удержание нормали и локальная глубина. Для ЭЧ это означает расплетание узла. Для чёрной дыры предельный режим описывается как переход объёмной динамики в поверхностную запись.

В этом смысле горизонт является не стенкой и не входом к точечной сингулярности, а поверхностью предельной фиксации:
\begin{equation}
\normal \to 0 \quad \Rightarrow \quad \text{объёмная память} \to \text{поверхностная запись}.
\end{equation}

Эта интерпретация согласуется с тем, что энтропия чёрной дыры пропорциональна площади горизонта, а не объёму. В языке модели это означает переход локальной памяти ЭЧ в глобальную форму записи.

Здесь важно явно отделить модель 01 от стандартной интерпретации испарения чёрных дыр. В полуклассическом описании излучение Хокинга рассматривается как термальный эффект, связанный с горизонтом, квантовыми полями и выбором доступных мод. Модель 01 принимает важность горизонта, энтропии и редуцированного описания, но не принимает буквальную картину чёрной дыры как термодинамического объекта, который постепенно испаряется и уносит информацию наружу тепловым потоком.

Рабочая гипотеза модели 01 иная:
\begin{equation}
\begin{aligned}
\text{чёрная дыра} &\neq \text{испаряющийся излучатель}, \\
\text{чёрная дыра} &\sim \text{поверхность фазовой записи}.
\end{aligned}
\end{equation}
В этом языке горизонт не теряет запись, а накапливает её. То, что в стандартной теории описывается как термальность горизонта, в модели 01 должно быть переинтерпретировано как ограничение доступа к фазовой информации и переход объёмной структуры в поверхностный код. Это сильное отличие модели и поэтому оно пока должно рассматриваться как гипотеза, требующая отдельной статьи и строгого сопоставления с расчётами Хокинга.

\section{Программа дальнейшей формализации}

На данном этапе модель 01 следует понимать не как завершённую физическую теорию, а как исследовательскую программу. Её сила состоит в том, что она предлагает единый язык для нескольких мотивов современной физики: поля, фазы, топологической устойчивости, наблюдательно-зависимого понятия частицы, памяти и горизонта. Но чтобы этот язык стал физической теорией, необходимо пройти несколько последовательных шагов.

Первый шаг --- определить пространство состояний поля 01. Пока поле записано символически как
\begin{equation}
\field=(0,1),
\end{equation}
но для строгой теории нужно указать, что именно является динамической переменной: скалярная фаза, многокомпонентная фаза, спинор, калибровочное поле или более общая геометрическая структура. Минимальная версия может начинаться с фазового поля $\phase(x)$ и амплитуды $A(x)$:
\begin{equation}
\psi(x)=A(x)e^{i\phase(x)}.
\end{equation}

Второй шаг --- задать функционал энергии или действие. Если частица является устойчивым узлом, то этот узел должен быть экстремумом некоторого функционала. В первом приближении можно искать выражение вида
\begin{equation}
E[\phase,A]=\int d^3x\,\left[\alpha |\nabla A|^2+\beta A^2|\nabla\phase|^2+V(A,\phase)+\gamma\,\mathcal{T}(\phase)\right],
\end{equation}
где $V$ задаёт локальную устойчивость, а $\mathcal{T}(\phase)$ обозначает возможный топологический вклад. Эта формула не является окончательной; её роль --- показать, какого типа математический объект требуется.

Третий шаг --- определить локальный скалярный/VEV-профиль $N$ не только словесно, но и геометрически. Возможный путь состоит в том, чтобы связать этот профиль с несовпадением двух фазовых ветвей, а затем добавить к нему интерпретацию удержания нормали:
\begin{equation}
N \sim \mathcal{N}\left[\phase_1-\phase_2\right], \qquad \normal \equiv \text{интерпретация}(N).
\end{equation}
Тогда подавление профиля будет означать не произвольное утверждение, а предел, в котором две фазовые ветви становятся неразличимыми:
\begin{equation}
\phase_1-\phase_2 \to 0 \quad \Rightarrow \quad N \to 0, \qquad \normal \to 0 \;\text{как интерпретационный предел}.
\end{equation}

Четвёртый шаг --- связать массу с энергией устойчивого решения. Если замкнутый узел является стационарной конфигурацией поля, то его масса должна выражаться через полную энергию этой конфигурации:
\begin{equation}
m_n c^2 = E_n,
\end{equation}
где индекс $n$ обозначает топологический или фазовый класс решения. Тогда формула \eqref{eq:mass_tension} станет не отдельной аналогией, а частным приближением более общего энергетического выражения.

Пятый шаг --- сопоставить внутренние числа модели с квантовыми числами. Нужно понять, может ли фазовый обход описывать заряд, спин, поколение, цветовой заряд и слабый изоспин. На этом этапе особенно важно не подменять Стандартную модель словами, а искать точные соответствия:
\begin{center}
\begin{tabular}{ll}
обход $\contour$ & топологический класс или число намотки, \\
ориентация обхода & возможный знак заряда, \\
скрученный обход & кандидат на спиновую структуру, \\
многокомпонентная фаза & кандидат на внутренние симметрии, \\
устойчивый минимум энергии & кандидат на массу частицы.
\end{tabular}
\end{center}

Шестой шаг --- вывести горизонт как предельное отображение объёмной структуры в поверхностную. Для этого нужно заменить словесную запись
\begin{equation}
\text{объёмная память} \to \text{поверхностная запись}
\end{equation}
на отображение между состояниями:
\begin{equation}
\mathcal{M}_{\mathrm{bulk}} \longrightarrow \mathcal{M}_{\mathrm{boundary}}.
\end{equation}
Если такое отображение сохраняет фазовую информацию, то модель сможет быть сопоставлена с голографическим принципом и энтропией Бекенштейна--Хокинга.

Седьмой шаг --- сформулировать возможные проверяемые отличия. Даже если модель сначала является интерпретационной, ей нужно постепенно искать места, где она может дать новые вопросы или предсказания: спектр устойчивых узлов, ограничения на массы, связь между фазовой памятью и энтропией, особенности ранних чёрных дыр, топологические аналоги в конденсированных средах.

Таким образом, ближайшая рабочая программа состоит не в том, чтобы сразу заменить существующую физику, а в том, чтобы построить мост:
\begin{equation}
\begin{aligned}
\text{образ модели 01}
&\longrightarrow \text{математическая структура} \\
&\longrightarrow \text{сопоставление с известной физикой}.
\end{aligned}
\end{equation}
Именно этот мост должен стать основой следующих статей.

\section{Что утверждается и что остаётся гипотезой}

В данной статье утверждается как рабочее ядро модели:
\begin{enumerate}
\item элементарная частица может быть описана как замкнутая фазовая структура;
\item масса связана с энергией удержания локального скалярного/VEV-профиля, интерпретируемого как удержание нормали;
\item фотон является открытым обходом без локального скалярного/VEV-профиля;
\item точечное материальное ядро заменяется динамическим ядром;
\item горизонт чёрной дыры является предельной поверхностной записью.
\end{enumerate}

Остаётся гипотезой и требует дальнейшей работы:
\begin{enumerate}
\item строгий вывод спектра частиц;
\item связь с калибровочными симметриями Стандартной модели;
\item количественная формула массы;
\item математическое определение памяти;
\item вывод горизонта и энтропии из фазовой динамики поля;
\item построение действия или энергетического функционала для поля 01;
\item поиск проверяемых следствий, отличающих модель от чистой интерпретации.
\end{enumerate}

\section{Заключение}

Модель 01 позволяет сформулировать элементарную частицу не как точечный объект, а как устойчивый замкнутый режим поля. В этом описании масса, спин, заряд и взаимодействия описываются через структуру обхода, локальный скалярный/VEV-профиль, его интерпретацию как удержания нормали и фазовую память. Такой подход не отменяет стандартную физику, но предлагает новый язык для её интерпретации: физический объект заменяется устойчивым ритмом, а точечное ядро --- динамической геометрией движения.

Первый результат этого языка состоит в устранении необходимости внутренней сингулярности: память и масса не помещаются в точку, а распределяются в фазовом обходе. На предельном уровне та же логика ведёт к горизонту как поверхности записи.

\end{document}